\section{Why Formation Matters For Intervention}

Once frameworks are understood as emergent historical structures, intervention
design changes. The central question is no longer only ``How do I resist the
bad state now?'' It becomes ``Which repeated local couplings are still
manufacturing the bad state, and which loop redesign will produce a different
macro-pattern later?''

This gives a more rigorous rationale for intervention:
\begin{enumerate}[nosep]
  \item identify the repeated cue-action-reward loop,
  \item identify which scale is carrying the regime,
  \item identify whether the regime is rigid, metastable, or fragile,
  \item redesign repeated conditions rather than relying only on heroic
        moment-level resistance.
\end{enumerate}

That is the business value of emergence language. It turns the manuscript from
a theory of resisting bad states into a theory of manufacturing better ones.
